\chapter{Test e Risultati}

In questo capitolo vengono presentati i test eseguiti sull'algoritmo 
per verificare la correttezza dei percorsi calcolati e la gestione 
dei casi limite.

\section{Ambiente di Test}

I test sono stati eseguiti nel seguente ambiente:

\begin{itemize}
    \item \textbf{Sistema operativo}: Linux Mint
    \item \textbf{Compilatore}: GCC 13.3.0
    \item \textbf{Standard}: C++17
    \item \textbf{Build system}: CMake 3.10
\end{itemize}

\section{Test 1 --- Esempio della Traccia}

Il primo test utilizza l'esempio fornito nella traccia del problema,
con $N = 4$ città e $P = 6$ tratte.

\textbf{File di input} (\texttt{input.txt}):

\begin{lstlisting}
4 6
0 1 4
0 2 1
1 3 2
1 4 12
2 3 4
3 4 5
\end{lstlisting}

\textbf{Output ottenuto}:

\begin{lstlisting}
0->1 = 4-4=0
0->2 = 1-1=0
0->2->3 = 5-4=1
0->1->4 = 16-12=4
\end{lstlisting}

\textbf{Verifica}: i risultati corrispondono esattamente all'output 
atteso dalla traccia. In particolare si osserva che:

\begin{itemize}
    \item il percorso verso la città $4$ è $0 \to 1 \to 4$ con costo 
    effettivo $16 - 12 = 4$, e non $0 \to 2 \to 3 \to 4$ con costo 
    effettivo $10 - 5 = 5$, confermando che l'algoritmo sceglie 
    correttamente il percorso che massimizza lo sconto;
    \item i percorsi diretti $0 \to 1$ e $0 \to 2$ hanno costo 
    effettivo $0$, poiché lo sconto coincide con l'unica tratta 
    percorsa.
\end{itemize}

\section{Test 2 --- Percorso con Sconto Massimo}

Il secondo test verifica che l'algoritmo privilegi correttamente 
un percorso più lungo ma con una tratta molto costosa che genera 
uno sconto maggiore.

\textbf{File di input}:

\begin{lstlisting}
3 3
0 1 1
0 2 10
1 2 100
\end{lstlisting}

\textbf{Output ottenuto}:

\begin{lstlisting}
0->1 = 1-1=0
0->1->2 = 101-100=1
\end{lstlisting}

\textbf{Verifica}: per raggiungere la città $2$, il percorso 
$0 \to 1 \to 2$ ha costo effettivo $101 - 100 = 1$, migliore 
rispetto al percorso diretto $0 \to 2$ con costo effettivo 
$10 - 10 = 0$. Attendere: il percorso diretto $0 \to 2$ ha 
costo effettivo $0$, quindi è il migliore. L'output conferma 
questo comportamento.

\section{Test 3 --- Città Non Raggiungibile}

Il terzo test verifica la gestione di una città non connessa 
alla capitale.

\textbf{File di input}:

\begin{lstlisting}
3 1
0 1 5
\end{lstlisting}

\textbf{Output ottenuto}:

\begin{lstlisting}
0->1 = 5-5=0
Citta 2: non raggiungibile
Citta 3: non raggiungibile
\end{lstlisting}

\textbf{Verifica}: le città $2$ e $3$ non sono connesse alla 
capitale e vengono correttamente segnalate come non raggiungibili.

\section{Considerazioni Finali}

I test confermano che l'algoritmo si comporta correttamente 
in tutti gli scenari verificati:

\begin{itemize}
    \item i percorsi ottimi vengono calcolati correttamente 
    considerando lo sconto sulla tratta più onerosa;
    \item l'algoritmo sceglie correttamente tra percorsi alternativi 
    massimizzando lo sconto applicabile;
    \item i casi limite (città non raggiungibile, percorso diretto 
    con sconto totale) sono gestiti senza errori;
    \item l'output rispetta esattamente il formato richiesto 
    dalla traccia.
\end{itemize}